% status: ZERO
\chapter{CPU and SIMD}\label{app:cpu}
\begin{ChapterIntent}
A CPU is both a numerical target and the controller of an accelerator
service. This reference connects loops to cache lines, vector lanes and
threads. Use it with Appendix A.5 and Chapters 3 and 7; instruction-specific
performance belongs to those chapters' measurement records.
\end{ChapterIntent}

\section{Elements are addressed; lines are transferred}
An array access names an element, but memory hardware commonly transfers
a cache line. Sequential access can use most of it. A large stride may
transfer many bytes for each useful element. On an illustrative machine
with 64-byte lines and four-byte elements, a line holds sixteen elements.
Reading one element per distinct line uses one sixteenth of its payload
before reuse. This is not the actual M1 line size.

Row-major traversal works well because neighbouring column indices name
neighbouring elements. Transposing a view changes strides without
necessarily moving data. A loop that ignores those strides can be wrong;
one that respects them can still be slow. Appendix A.4 separates these
failures. Neither a contiguous allocation nor a matrix-shaped API proves
that a chosen traversal has locality.

\section{Vector lanes and the tail}
Single instruction, multiple data (SIMD) applies an operation across
several lanes. It helps when data can feed those lanes and the compiler
can establish safe accesses. Branching, aliasing, strides and reductions
can complicate that proof. An intrinsic makes an operation explicit but
does not make the surrounding memory access efficient.

For an illustrative eight-lane vector processing nineteen elements, two
full vectors leave three elements. A scalar cleanup or correctly masked
operation handles that tail. Loading eight elements from a three-element
region without a valid memory contract is not an optimization. Test
lengths just below and above vector boundaries, and test misaligned views
if the interface allows them.

NEON, AVX and matrix extensions differ in width, operations and dispatch
requirements. Check the actual target and runtime support rather than
inferring support from the operating system. A portable scalar reference
is valuable even when the optimized path is architecture-specific.

\section{Reduction order is part of the result}
Several independent accumulators can expose instruction-level parallelism.
Combining them changes summation order. Real-number algebra allows this;
floating-point arithmetic need not produce identical bits. A random seed
does not repair an altered arithmetic result. Appendix~\ref{app:numerics}
explains the numerical contract.

First compare the output against an independent reference. Then inspect
vectorization reports or generated instructions. Finally measure the
intended data size and thread count. A tiny array in cache answers a
different question from a checkpoint larger than the last-level cache.

\section{Threads share resources}
Threads may share caches and a memory interface. Once aggregate traffic
saturates the interface, more threads can add overhead without increasing
throughput. False sharing occurs when threads write different values in
the same coherence line: logically independent updates still trigger
ownership transfers. Separate per-thread accumulators and combine them
deliberately.

Record affinity, the cores actually used, and nested library thread counts.
Giving every nested layer all available cores can oversubscribe the
machine. On a NUMA system, page placement and the accessing core determine
which memory path is used. Establish locality before calling one result
the machine's memory bandwidth.

\section{The CPU also runs the service}
A CPU path can win for tiny operations whose accelerator submission costs
more than their execution. The same CPU may tokenize, enforce grammars,
schedule batches and service streams. A faster isolated matrix kernel can
therefore worsen a server if it occupies control-plane cores.

Predict the shared resource, measure a thread-count sweep, and retain a
one-thread baseline. Never multiply a vector-width ratio by a core-count
ratio and present the product as measured speedup. The roofline and
launch-overhead chapters demonstrate how to test that hypothesis.
